% !TEX root = ../../../proposal.tex

\section{Diffusion Models}
\label{sec:diffusion_models}

Diffusion models \citep{pmlr-v37-sohl-dickstein15} are latent-variable generative models of the form
\(
p_\theta(x_0)
=
\int p_\theta(x_{0:T})\, \mathrm{d}x_{1:T},
\)
where \(x_{1:T}\) are latents of the same dimensionality as the observed data
\(x_0\). Following \citet{NEURIPS2020_4c5bcfec}, the model consists of a fixed
\emph{forward diffusion process} that gradually adds noise to data, and a learned
\emph{reverse process} that aims to invert this corruption.

\begin{figure}[H]
    \centering
    \includegraphics[width=0.85\textwidth]{chapters/background/diffusion_models/media/DDPM.png}
    \caption{Directed graphical model structure of DDPMs \citep{NEURIPS2020_4c5bcfec}.}
    \label{fig:ddpm_graphical_model}
\end{figure}

The joint distribution of all variables is written as a Markov chain:
\begin{equation}
    p_\theta(x_{0:T})
    =
    p(x_T)
    \prod_{t=1}^{T} p_\theta(x_{t-1}\mid x_t),
    \qquad
    p(x_T)=\mathcal{N}(0,I),
    \label{eq:ddpm_joint}
\end{equation}
where \(p_\theta(x_{t-1}\mid x_t)\) is the reverse process. The forward diffusion
process is defined as a separate Markov chain:
\begin{equation}
    q(x_{1:T}\mid x_0)
    =
    \prod_{t=1}^T
    q(x_t\mid x_{t-1}),
    \qquad
    q(x_t\mid x_{t-1})
    =
    \mathcal{N}\!\left(
        \sqrt{1-\beta_t}\,x_{t-1},\;
        \beta_t I
    \right),
    \label{eq:ddpm_forward}
\end{equation}
where \(\{\beta_t\}_{t=1}^T\) is the variance schedule that \emph{can} be learned by reparameterization,
or fixed to predetermined constants (e.g., a linear schedule from $10^{-4}$ to $10^{-2}$) such that \(\beta_t\) typically increases slowly with \(t\).
This ensures that the forward chain remains
analytically tractable and preserves the same functional form across timesteps while gradually
driving \(x_T\) toward an isotropic Gaussian.

The authors \citep{NEURIPS2020_4c5bcfec} choose \(T\) large (e.g.\ \(T=1000\)) and let \(\beta_t\) increase linearly or according
to a smooth curve from a very small value to a modest value. This ensures that each step only
slightly degrades the signal, so the forward chain remains close to a diffusion process with
well-behaved Gaussian conditionals, and by the final step \(x_T\) is nearly standard
Gaussian.

%%%%%%%%%%%%%%%%%%%%%%%%%%%%%%%%%%%%%%%%%%%%%%%%%%%%%%%%%%%%%%%%%%%%%%%%%%%%%%%
\subparagraph{Forward Diffusion Process.}
%%%%%%%%%%%%%%%%%%%%%%%%%%%%%%%%%%%%%%%%%%%%%%%%%%%%%%%%%%%%%%%%%%%%%%%%%%%%%%%

Because each transition in \eqref{eq:ddpm_forward} is linear Gaussian, the forward chain admits a
closed-form marginal at any timestep and repeatedly composing the Gaussian transitions yields:
\begin{equation}
    q(x_t\mid x_0)
    =
    \mathcal{N}\!\left(
        \sqrt{\bar\alpha_t}\,x_0,\;
        (1-\bar\alpha_t)I
    \right)
    \text{ where} \quad
    \alpha_t := 1-\beta_t,
    \qquad
    \bar\alpha_t := \prod_{s=1}^{t}\alpha_s.
    \label{eq:ddpm_forward_marginal}
\end{equation}
We can define \(x_t\) through an explicit reparameterization
\begin{align}
    x_t &= \sqrt{\alpha_t}\,x_{t-1}
    +
    \sqrt{1-\alpha_t}\,\varepsilon_{t-1} \\
    &= \sqrt{\alpha_t \alpha_{t-1}}\,x_{t-2}
    +
    \sqrt{\alpha_t(1-\alpha_{t-1})}\,\varepsilon_{t-2}
    +
    \sqrt{1-\alpha_t}\,\varepsilon_{t-1} \\
    &= \dots \\
    &=
    \sqrt{\bar\alpha_t}\,x_0
    +
    \sqrt{1-\bar\alpha_t}\,\varepsilon,
    \qquad
    \varepsilon\sim\mathcal{N}(0,I).
    \label{eq:reparameterization_of_xt}
\end{align}
where the state at timestep \(t\) is a combination of the original data and Gaussian noise,
with the signal coefficient \(\sqrt{\bar\alpha_t}\) gradually shrinking as \(t\) increases.

%%%%%%%%%%%%%%%%%%%%%%%%%%%%%%%%%%%%%%%%%%%%%%%%%%%%%%%%%%%%%%%%%%%%%%%%%%%%%%%
\subparagraph{Reverse Denoising Process.}
%%%%%%%%%%%%%%%%%%%%%%%%%%%%%%%%%%%%%%%%%%%%%%%%%%%%%%%%%%%%%%%%%%%%%%%%%%%%%%%

The reverse process attempts to invert the forward diffusion by learning Gaussian transitions of the
form
\begin{equation}
    p_\theta(x_{t-1}\mid x_t)
    =
    \mathcal{N}\!\left(
        \mu_\theta(x_t,t),\;
        \Sigma_\theta(x_t,t)
    \right).
    \label{eq:ddpm_reverse}
\end{equation}
Here, \(\theta\) denotes a single shared neural-network parameterization; the dependence on \(t\) is
provided through a timestep embedding. Because the forward process \(q(x_{1:T}\mid x_0)\) is fixed and Gaussian, the
associated conditional distributions \(q(x_{t-1}\mid x_t,x_0)\) admit
closed-form expressions. These conditionals arise directly from standard Gaussian identities:
they describe how the forward noising process relates adjacent latent variables when conditioned
on the original clean sample \(x_0\). The DDPM model parameterizes
\(p_\theta(x_{t-1}\mid x_t)\) to approximate these analytically derived
Gaussian conditionals.
It has the Gaussian form
\begin{equation}
    q(x_{t-1}\mid x_t,x_0)
    =
    \mathcal{N}\!\left(
        \tilde{\mu}_t(x_t,x_0),\;
        \tilde{\beta}_t I
    \right),
    \label{eq:true_reverse_posterior}
\end{equation}
with
\begin{align}
    \tilde{\mu}_t(x_t,x_0)
    &=
    \frac{\sqrt{\bar\alpha_{t-1}}\beta_t}{1-\bar\alpha_t}\,x_0
    +
    \frac{\sqrt{1-\beta_t}(1-\bar\alpha_{t-1})}{1-\bar\alpha_t}\,x_t,
    \\
    \tilde{\beta}_t
    &=
    \frac{1-\bar\alpha_{t-1}}{1-\bar\alpha_t}\,\beta_t.
\end{align}
The model therefore learns \(p_\theta(x_{t-1}\mid x_t)\) so that its mean matches
the analytically derived forward-process posterior. For stability, DDPM fixes the reverse-process
variance to one of two experimentally evaluated choices:
\[
    \sigma_t^2 = \beta_t
    \qquad\text{or}\qquad
    \sigma_t^2 = \tilde{\beta}_t = \frac{1-\bar{\alpha}_{t-1}}{1-\bar{\alpha}_t}\,\beta_t,
\]
as reported by \citet{NEURIPS2020_4c5bcfec}. Both forms preserve the Gaussian functional structure
needed for closed-form KL divergences in the ELBO, while yielding similar empirical performance.

%%%%%%%%%%%%%%%%%%%%%%%%%%%%%%%%%%%%%%%%%%%%%%%%%%%%%%%%%%%%%%%%%%%%%%%%%%%%%%%
\subparagraph{Training Objective.}
%%%%%%%%%%%%%%%%%%%%%%%%%%%%%%%%%%%%%%%%%%%%%%%%%%%%%%%%%%%%%%%%%%%%%%%%%%%%%%%

Training minimizes the negative log-likelihood \(-\log p_\theta(x_0)\). Since
\(x_{1:T}\) is intractable to marginalize exactly, one optimizes a standard variational
lower bound:
\begin{equation}
    \mathbb{E}_q[-\log p_\theta(x_0)]
    \;\le\;
    \mathcal{L}(\theta)
    =
    \mathbb{E}_q\!\left[
        -\log\frac{p_\theta(x_{0:T})}{q(x_{1:T}\mid x_0)}
    \right]
    =
    \mathbb{E}_q\!\left[
        -\log p(x_T)
        -
        \sum_{t=1}^T
        \log\frac{p_\theta(x_{t-1}\mid x_t)}
        {q(x_t\mid x_{t-1},x_0)}
    \right]
    =:
    L(\theta),
    \label{eq:elbo_ddpm_raw}
\end{equation}
This expression can be rewritten as a sum of KL terms \citep{weng2021diffusion}:
\begin{equation}
    \mathcal{L}(\theta)
    =
    \mathbb{E}_q\Big[
        \underbrace{
            D_{\mathrm{KL}}\bigl(q(x_T\mid x_0)\,\|\,p(x_T)\bigr)
        }_{L_T}
        +
        \sum_{t=2}^{T}
        \underbrace{
            D_{\mathrm{KL}}\bigl(
                q(x_{t-1}\mid x_t,x_0)
                \;\|\;
                p_\theta(x_{t-1}\mid x_t)
            \bigr)
        }_{L_{t-1}}
        -
        \underbrace{\log p_\theta(x_0\mid x_1)}_{L_0}
    \Big].
    \label{eq:elbo_ddpm_kl}
\end{equation}
The term \(L_T\) encourages the forward terminal distribution to match the prior. Each \(L_{t-1}\)
penalizes discrepancies between the true reverse posterior and the learned reverse conditional at
that timestep. The final term \(L_0\) plays the role of a reconstruction likelihood for
\(x_0\). Because all of these terms compare Gaussian distributions with known parameters,
they admit closed-form expressions, allowing efficient estimation with low-variance gradients.

%%%%%%%%%%%%%%%%%%%%%%%%%%%%%%%%%%%%%%%%%%%%%%%%%%%%%%%%%%%%%%%%%%%%%%%%%%%%%%%
%%%%%%%%%%%%%%%%%%%%%%%%%%%%%%%%%%%%%%%%%%%%%%%%%%%%%%%%%%%%%%%%%%%%%%%%%%%%%%%
\subparagraph{Simplified Noise-Prediction Objective.}
%%%%%%%%%%%%%%%%%%%%%%%%%%%%%%%%%%%%%%%%%%%%%%%%%%%%%%%%%%%%%%%%%%%%%%%%%%%%%%%

The variational bound in \eqref{eq:elbo_ddpm_kl} has the generic form
\(
    L_{\mathrm{VLB}} = L_T + \sum_{t=2}^T L_{t-1} + L_0,
\)
where \(L_T\) and \(L_0\) do not depend on the parametrization of the reverse mean.
The terms \(L_{t-1}\) are KL divergences between Gaussians and therefore can be written in closed
form. Assuming an isotropic reverse covariance
\(\Sigma_\theta(x_t,t)=\sigma_t^2 I\), each such term reduces to
a mean-squared error between the true posterior mean \(\tilde{\mu}_t\) and the model
mean \(\mu_\theta\):
\begin{equation}
    L_{t-1}
    =
    \mathbb{E}_q\!\left[
        \frac{1}{2\sigma_t^2}
        \bigl\|
            \tilde{\mu}_t(x_t,x_0)
            -
            \mu_\theta(x_t,t)
        \bigr\|^2
    \right]
    + C_t,
    \label{eq:L_t_mse_form}
\end{equation}
where \(C_t\) does not depend on \(\theta\).

Using the closed-form forward posterior
\(
    q(x_{t-1}\mid x_t,x_0)
    =
    \mathcal{N}(\tilde{\mu}_t,\tilde{\beta}_t I)
\)
and the forward reparameterization
\(
    x_t = \sqrt{\bar{\alpha}_t}\,x_0 +
    \sqrt{1-\bar{\alpha}_t}\,\varepsilon,\;
    \varepsilon\sim\mathcal{N}(0,I),
\)
the posterior mean can be expressed as
\begin{equation}
    \tilde{\mu}_t(x_t,x_0)
    =
    \frac{1}{\sqrt{\alpha_t}}
    \left(
        x_t
        -
        \frac{1-\alpha_t}{\sqrt{1-\bar{\alpha}_t}}\,
        \varepsilon
    \right),
\end{equation}
Since \(x_t\) is available as model
input, it is convenient to parametrize the reverse mean so that the network predicts the noise
directly:
\begin{equation}
    \mu_\theta(x_t,t)
    =
    \frac{1}{\sqrt{\alpha_t}}
    \left(
        x_t
        -
        \frac{1-\alpha_t}{\sqrt{1-\bar{\alpha}_t}}\,
        \varepsilon_\theta(x_t,t)
    \right).
    \label{eq:reverse_mean_eps_param}
\end{equation}
Substituting \(\tilde{\mu}_t\) and \(\mu_\theta\) into
\eqref{eq:L_t_mse_form} and simplifying gives
\begin{equation}
    L_{t-1}
    =
    \mathbb{E}_{x_0,\varepsilon}
    \left[
        \frac{(1-\alpha_t)^2}{2\sigma_t^2\,\alpha_t(1-\bar{\alpha}_t)}
        \,
        \bigl\|
            \varepsilon
            -
            \varepsilon_\theta(x_t,t)
        \bigr\|^2
    \right]
    + C_t,
\end{equation}
so each KL term is a weighted mean-squared error between the true noise
\(\varepsilon\) and the predicted noise \(\varepsilon_\theta(x_t,t)\).

Dropping the timestep-dependent weight yields the commonly used simplified objective
\begin{equation}
    L_{\mathrm{simple}}(\theta)
    =
    \mathbb{E}_{t,x_0,\varepsilon}
    \bigl[
        \|
            \varepsilon
            -
            \varepsilon_\theta(x_t,t)
        \|^2
    \bigr],
    \qquad
    x_t
    =
    \sqrt{\bar{\alpha}_t}\,x_0
    +
    \sqrt{1-\bar{\alpha}_t}\,\varepsilon,
    \label{eq:ddpm_simple_loss}
\end{equation}
which is equivalent to a reweighted form of the ELBO but is simpler to implement and has been
observed to yield improved training stability and sample quality in practice.

%%%%%%%%%%%%%%%%%%%%%%%%%%%%%%%%%%%%%%%%%%%%%%%%%%%%%%%%%%%%%%%%%%%%%%%%%%%%%%%
\subparagraph{Training and Sampling Algorithms.}

The training and sampling procedures are summarized in Figure~\ref{fig:algorithms}. Training draws
a data sample \(x_0\), chooses a timestep \(t\) uniformly from \(\{1,\dots,T\}\), samples
\(\varepsilon\sim\mathcal{N}(0,I)\), constructs \(x_t\) using the forward
marginal \eqref{eq:ddpm_forward_marginal}, and takes a gradient step on the objective
\eqref{eq:ddpm_simple_loss}. Sampling inverts this process: starting from
\(x_T\sim\mathcal{N}(0,I)\), the model repeatedly applies the reverse
transition implied by \eqref{eq:reverse_mean_eps_param},
\[
    x_{t-1}
    =
    \frac{1}{\sqrt{\alpha_t}}
    \left(
        x_t
        -
        \frac{\beta_t}{\sqrt{1-\bar\alpha_t}}\,
        \varepsilon_\theta(x_t,t)
    \right)
    +
    \sigma_t z,
    \qquad
    z\sim\mathcal{N}(0,I),
\]
for \(t = T,T-1,\dots,1\), where \(\sigma_t^2\) is chosen consistently with the forward variance
schedule. This iterative denoising procedure gradually transforms white noise into a structured
sample \(x_0\).

%%%%%%%%%%%%%%%%%%%%%%%%%%%%%%%%%%%%%%%%%%%%%%%%%%%%%%%%%%%%%%%%%%%%%%%%%%%%%%%
\subparagraph{Architecture.}

The reverse-process mean or noise predictor is implemented using a U-Net architecture with residual
blocks, skip connections, and multi-scale feature hierarchies \citep{NEURIPS2020_4c5bcfec}. The
timestep \(t\) is encoded by a sinusoidal positional embedding and injected into intermediate layers,
allowing the network to adapt its predictions to the current noise level.
The network outputs
\(\varepsilon_\theta(x_t,t)\), which is converted to the reverse mean via
Equation~\eqref{eq:reverse_mean_eps_param}, and the variance is typically fixed as
\(\Sigma_\theta(x_t,t)=\sigma_t^2 I\) to keep training stable.

\subsection{Score-Based Formulation}

Song et al.~\citep{song2021scorebased} showed that DDPMs can be understood in continuous time
through stochastic differential equations. In this view, discrete-time diffusion models and
score-based generative models are two parameterizations of the same underlying denoising process.

\begin{figure}[H]
    \centering
    \includegraphics[width=0.65\textwidth]{chapters/background/diffusion_models/media/score.jpg}
    \caption{\textbf{Forward and reverse SDEs.} In continuous time, the forward process corrupts
    data with drift and diffusion terms, while the reverse process uses the score function to
    steer noisy samples back toward the data manifold \citep{song2021scorebased}.}
    \label{fig:sde_score_diagram}
\end{figure}

\paragraph{From the Markov Chain to the Continuous Limit.}
The discrete forward chain can be viewed as a time discretization of a continuous corruption
process. For variance-preserving diffusion, taking the small-step limit yields
\begin{equation}
    \mathrm{d}x
    =
    -\tfrac{1}{2}\beta(t)\,x\,\mathrm{d}t
    +
    \sqrt{\beta(t)}\,\mathrm{d}W(t),
\end{equation}
where $W(t)$ is a Wiener process. More generally, the forward corruption process is written as
\begin{equation}
    \mathrm{d}x = f(x,t)\,\mathrm{d}t + g(t)\,\mathrm{d}w,
\end{equation}
with drift \(f(x,t)\) and diffusion strength \(g(t)\).

\paragraph{Reverse SDE.}
The corresponding reverse-time SDE has the form
\begin{equation}
    \mathrm{d}x
    =
    \left[
        f(x,t) - g(t)^2 \nabla_x \log p_t(x)
    \right]\mathrm{d}t
    +
    g(t)\,\mathrm{d}\bar w,
\end{equation}
where \(\nabla_x \log p_t(x)\) is the \emph{score function}, namely the gradient of the
log-density at noise level \(t\). Intuitively, the score points in the direction that makes
the current sample more likely under the time-dependent data distribution. The reverse SDE is
the continuous-time counterpart of the DDPM reverse chain where reverse denoising is driven by a vector field
that points back toward the data manifold.

In the DDPM parameterization, the model outputs \(\varepsilon_\theta(x_t,t)\), an estimate
of the Gaussian noise used to corrupt \(x_0\) into \(x_t\). Because \(x_t\) is a linear
combination of signal and noise, this noise estimate can be converted directly into a score
estimate:
\begin{equation}
    s_\theta(x_t, t) = -\frac{\varepsilon_\theta(x_t, t)}{\sqrt{1-\bar\alpha_t}}.
    \label{eq:score_noise}
\end{equation}
DDPMs can therefore be viewed as score-based generative models written in a different
parameterization. 
The connection is important for
Phase~1 because the compositional operators studied there combine score functions from
separately conditioned model evaluations, rather than modifying model weights or model
architecture.


\subsection{Conditional Generation and Classifier-Free Guidance}

This section distinguishes between \textbf{unconditional} and \textbf{conditional} diffusion models by examining their objectives, formulations, and applications.

While the above formulation defines the foundation of \emph{unconditional diffusion models}, practical applications often demand controlled or semantically guided generation such as synthesizing specific classes, attributes, or text-aligned imagery.
This introduces the need for \emph{conditional diffusion models}, where additional information \(y\) (e.g., labels, text embeddings, or medical attributes) guides the reverse process.
Recent extensions, such as \emph{classifier} \citep{NEURIPS2021_49ad23d1} and \emph{classifier-free guidance} \citep{ho2022classifier}, further refine this conditioning by explicitly or implicitly modulating the score function to balance fidelity and diversity during sampling.

\paragraph{Unconditional Diffusion Models.}
An \textit{unconditional diffusion model} learns to generate samples from the marginal data distribution \(p(x)\) without using auxiliary information or labels. During training, the model observes only samples \(x_0\) drawn from \(p_{\text{data}}(x)\) and is optimized to predict the noise added at each diffusion step:
\begin{equation}
\mathcal{L}_{\text{uncond}} = \mathbb{E}_{x_0, t, \varepsilon}\!\left[\|\varepsilon - \varepsilon_\theta(x_t, t)\|^2\right].
\end{equation}
The training signal is generated from the data itself rather than supplied by external annotation, so the objective is self-supervised. At sampling time, the model starts from Gaussian noise \(x_T \sim \mathcal{N}(0, I)\) and denoises it step by step to produce a realistic sample. In this sense, unconditional diffusion models learn \(p(x)\) directly, without any conditioning variable.

\paragraph{Conditional Diffusion Models.}
A \textit{conditional diffusion model} extends the unconditional setting by modeling
\(p(x \mid y)\), where \(y\) provides additional information such as a class label, text
prompt, or spatial map. The forward noising process is unchanged. Conditioning enters only
through the reverse model, which is now written as
\[
p_\theta(x_{t-1}\mid x_t,y).
\]
In the DDPM parameterization, this is implemented by conditioning the noise predictor
\(\varepsilon_\theta(x_t,t,y)\). Training therefore uses the objective
\begin{equation}
\mathcal{L}_{\text{cond}}
=
\mathbb{E}_{x_0,t,\varepsilon,y}
\left[
\|\varepsilon-\varepsilon_\theta(x_t,t,y)\|^2
\right].
\end{equation}
By Equation~\ref{eq:score_noise}, this is equivalent to learning the conditional score of
the noisy data distribution. At sampling time, the condition \(y\) steers the denoising
trajectory toward samples consistent with the desired class, prompt, or attribute. In this
way, conditional diffusion turns unconditional generation into controllable generation.

\paragraph{Classifier Guidance.}
Classifier guidance \citep{NEURIPS2021_49ad23d1} provides a post-hoc way to steer an
unconditional diffusion model toward a target condition \(y\). The key observation is that
a classifier gives access to \(p_\phi(y\mid x_t)\), and therefore to the gradient
\[
\nabla_{x_t}\log p_\phi(y\mid x_t).
\]
Applying Bayes' rule gives
\[
p_t(x_t\mid y)=\frac{p_\phi(y\mid x_t)\,p_t(x_t)}{p(y)},
\]
so
\[
\log p_t(x_t\mid y)
=
\log p_\phi(y\mid x_t)+\log p_t(x_t)-\log p(y),
\]
and differentiation with respect to \(x_t\) yields
\begin{equation}
\nabla_{x_t}\log p_t(x_t\mid y)
=
\nabla_{x_t}\log p_t(x_t)
+
\nabla_{x_t}\log p_\phi(y\mid x_t),
\label{eq:classifier_guidance_score}
\end{equation}
since \(p(y)\) does not depend on \(x_t\). The conditional score can therefore be written
as the sum of the unconditional score and a classifier-derived conditioning term. In
practice, \citet{NEURIPS2021_49ad23d1} scale this conditioning term by a guidance weight
\(\gamma \ge 0\):
\begin{equation}
\tilde{s}(x_t,t,y)
=
s_\theta(x_t,t)
+
\gamma\,\nabla_{x_t}\log p_\phi(y\mid x_t).
\label{eq:classifier_guidance_scaled}
\end{equation}
Larger values of \(\gamma\) strengthen adherence to the target condition, often at the cost
of reduced sample diversity.

\paragraph{Classifier-Free Guidance.}
Classifier-free guidance (CFG) \citep{ho2022classifier} achieves a similar effect without
training a separate classifier. A single diffusion model is trained with condition dropout
so that it learns both the conditional and unconditional predictors
\[
\varepsilon_\theta(x_t,t,y)
\qquad\text{and}\qquad
\varepsilon_\theta(x_t,t,\emptyset),
\]
where \(\emptyset\) denotes the null condition. At sampling time, these two predictions are
combined as
\begin{equation}
\tilde{\varepsilon}_\theta(x_t,t,y)
=
(1+w)\,\varepsilon_\theta(x_t,t,y)
-
w\,\varepsilon_\theta(x_t,t,\emptyset),
\label{eq:guided_eps_cfg}
\end{equation}
with guidance weight \(w \ge 0\). When \(w=0\), this reduces to ordinary conditional
generation. Larger values of \(w\) strengthen the conditional signal while reducing
diversity. This formulation has become the standard conditioning mechanism in modern
diffusion models. In the proposed study, it plays two concrete roles. 
This is relevant throughout the proposal because classifier-free guidance provides the basic conditioning mechanism used for standard prompting and for composition. 
In Phase~1, it underlies both the single-prompt baseline and the compositional operator of \citet{liu2022compositional}. 
In Phases~2 and~3, the same conditional prediction framework is used in the controlled synthetic and clinical composition settings.